\subsection*{从同一份初态比较三种写法}

第四章已经比较固定赋值与减法。这里各条余额历史继续从独立的 $b=10000$ 分开始：连续赋值 7000 保持 7000，连续减 3000 则到 4000；把另一项合法的减 1000 插在两次旧赋值之间，就得到 $10000\to7000\to6000\to7000$。后一个 7000 覆盖了中间支出，说明余额观察上的连续幂等，尚未处理交错更新。下面换成能保留中间新增信息的集合表达。

再看与购买有关的标签。设 O1 的标签集合 L 初始为空，t 表示加入 \code{gift}，u 表示加入 \code{priority}。t 把 L 改为 $L\cup\{\mathrm{gift}\}$。从空集合连续执行两次 t，结果都是 $\{\mathrm{gift}\}$。执行 t、u、t，最终得到 $\{\mathrm{gift},\mathrm{priority}\}$，中间加入的另一个标签得以保留。

这里的 $\cup$ 表示集合并：结果包含任一输入集合中的元素，相同元素只出现一次。\code{gift} 是标签值，也成为这一业务事实的身份。如果产品规则就是这张订单具有礼物标签，两次相同设置可以汇到同一事实。若规则改成每点击一次就赠送一份包装，两次点击可能各有独立意义，标签集合就缺少表达份数的位置。

两次合法购买即使商品相同，也必须分别存在。两次把 O1 标为礼物，在本例规则下则可以合为一个事实。购买入口用操作身份区别两项意图，标签对象直接按标签值保存成员资格。能否合并这两次输入，先由业务要保留的区别决定。

\subsection*{进度值背后需要哪些已确认事实}

第十章的消费事务 Q 已经按 I 识别 EV1 的重复接收，按业务屏障 B 仲裁 O1 的初次履约，取得真实编号为 501 的 F1，再把 O1 从 \code{unfulfilled} 推进到 \code{fulfillment\_accepted}。建立 F1 与实际发货分开：后续发送者从持久任务取得原 \code{dispatch} 身份，向外部履约系统继续执行。

本章在这条链上增加一个履约查询投影。投影是为查询整理的状态，它依次展示未受理、已受理、已发货、已送达。对这条单向正常路径，分别赋等级 0、1、2、3，来件到达时取当前等级与来件等级中较大的一个。

共同初态取 O1 的投影等级为 0。先收到已发货 2，再迟到收到已受理 1，结果依次为 2、2。重复收到 2，等级也保持 2。这样的表达容许通知重复与重排，条件是通知所指事实真实、属于同一 O1/F1，并且高等级所表示的阶段包含低等级要求的前驱事实。

如果某个未经核验的来件只写 \code{delivered}，给它数字 3 再取最大，就可能直接显示一件从未交给承运商的包裹已经送达。接收端先核对业务证据，才能把来件转成有依据的等级。第~\ref{section:13-4}~节会把这些依据展开为事实集合，再从集合计算显示阶段。

条件迁移还有另一类含义。例如扣款只在余额足够时执行，取消只在尚未发货时受理。前置条件会随其他操作改变。第一次判断成功，第二次在不同状态下重新判断，结果就可能变化。第五章保存首次确定的拒绝，第八章共同保存已受理的 P/C/T，让同一项操作在后来的状态中找到原结论或原恢复任务。

\CThirteenFigure{ch13-original02}{标签并与进度取最大合并已形成的信息，付款和发货从各自入口产生业务事实}{fig:ch13-objects}

购买、标签、投影与外部动作可以出现在同一应用中。平台通知 PEV 带来经核验的 C1/X1 证据，R 据此建立本地已付款事件 EV1，履约库再按 I/B 建立 F1，查询投影接收 F1 的进度。标签的集合并发生在这条链中的标签对象上。重投 EV1 时，消费者沿原 I/B 找回实际 F1；重发履约任务时，发送者沿原 \code{dispatch} 身份取得外部结局，随后把已确认的进度交给投影。
\FloatBarrier
